\section{Conclusion and limitations}

\paragraph{(1) Complexity.}
XGBoost [3] reaches test AUC 0.7887 / AP 0.5784 against 0.7636 / 0.5234 for the logistic
scorecard and 0.7706 for the KNN benchmark; the field spans 0.05 AUC, the ensembles are tied
(XGBoost over the forest, 0.0020 [$-0.0010$, $+0.0051$]) and nested cross-validation confirms
the frozen choices: boosting gains only what step effects in arrears are worth.

\paragraph{(2) Operating point.}
The cut-off is worth more than the model at every cost ratio from 2:1 to 20:1: at 5:1,
moving the logit from 0.5 to 0.15 saves 1,937 [1,604, 2,218] of 7,063 test cost units,
switching to XGBoost a further 145 [$-$18, 299], an interval that includes zero.

\paragraph{(3) Imbalance.}
Nothing improves ranking: weights, SMOTE and under-sampling inside the folds leave the
logit within 0.0012 AUC and cost the forest 0.0076; re-weighting moves the intercept
($+1.24$), buying recall with calibration (Brier 0.138 to 0.190) that Platt scaling restores
(0.139). Resampling before the split fabricates 0.914.

\paragraph{(4) Inference.}
Recent and repeated arrears dominate --- mean $|$SHAP$|$ 0.298, 0.281 and 0.230 log-odds for
PAY\_1, MAX\_DELAY and N\_MONTHS\_DELAYED, all pushing towards default --- and SHAP tracks
out-of-bag permutation (Spearman 0.77) while impurity importance tracks neither (0.31).
Friedman's $H = 0.061$ for PAY\_1 $\times$ LIMIT\_BAL refutes the proposal's interaction
hypothesis; monotone constraints cost 0.0016 AUC and remove every violation; the K-means
segments add nothing to the logit and are descriptive only.

\paragraph{Limitations.}
One split of one 2005 Taiwanese portfolio: the figures are a point estimate for one economy
and vintage, and the bootstrap prices test-row noise only. Six
months of history caps what any model can learn, the likeliest reason no family passes 0.79
AUC. SHAP and permutation describe the fitted model, not causation, and the
PAY\_RATIO\_MEAN zero bucket conflates 936 imputed non-borrowers (37.0\% default in
training) with genuine non-payers. The test set was read by 46 frozen scorings and by the
reporting-only bootstrap and cost curves: a report card, not a fresh holdout.

\paragraph{Future work.}
Repeated splits and a fresh holdout for honest intervals, with an indicator separating the
imputed zero bucket from genuine non-payment. Bureau, income and transaction data over a
longer window, measured loss-given-default in place of 5:1, and isotonic in place of Platt
recalibration.

\section*{References}

\footnotesize\setlength{\parskip}{2pt}

[1] Breiman, L.\ (2001) Random forests. {\it Machine Learning} {\bf 45}(1):5--32.

[2] Chawla, N.\,V., Bowyer, K.\,W., Hall, L.\,O.\ \& Kegelmeyer, W.\,P.\ (2002) SMOTE:
synthetic minority over-sampling technique. {\it J.\ Artificial Intelligence Research}
{\bf 16}:321--357.

[3] Chen, T.\ \& Guestrin, C.\ (2016) XGBoost: a scalable tree boosting system. {\it Proc.\ 22nd
ACM SIGKDD}, pp.\ 785--794.

[4] Friedman, J.\,H.\ (2001) Greedy function approximation: a gradient boosting machine.
{\it Annals of Statistics} {\bf 29}(5):1189--1232.

[5] Hofmann, H.\ (1994) {\it Statlog (German Credit Data)}. UCI Machine Learning Repository,
dataset 144.

[6] James, G., Witten, D., Hastie, T.\ \& Tibshirani, R.\ (2021) {\it An Introduction to
Statistical Learning}, 2nd ed. Springer.

[7] Lundberg, S.\,M.\ \& Lee, S.-I.\ (2017) A unified approach to interpreting model
predictions. {\it Advances in Neural Information Processing Systems} {\bf 30}:4765--4774.

[8] Nadeau, C.\ \& Bengio, Y.\ (2003) Inference for the generalization error. {\it Machine
Learning} {\bf 52}(3):239--281.

[9] Platt, J.\,C.\ (1999) Probabilistic outputs for support vector machines and comparisons
to regularized likelihood methods. In {\it Advances in Large Margin Classifiers}, MIT Press,
pp.\ 61--74.

[10] Yeh, I.-C.\ \& Lien, C.-H.\ (2009) The comparisons of data mining techniques for the
predictive accuracy of probability of default of credit card clients. {\it Expert Systems
with Applications} {\bf 36}(2):2473--2480.
